\subsection{Integer widths, bounded row padding, and mixed contractions}
\label{sec:mixed-width-rows}

The batched residual construction uses a fixed finite list of child sizes,
not only powers of one radix.  The following implementation supplies the
corresponding integer-width recursion.  All constants in this subsection
are fixed finite circuit constants, independent of the input length.

\paragraph{A general moment condition.}
Let $m,W\ge2$ be fixed.  At one circuit invocation, let the list
$t_1,\ldots,t_J$ contain integers in $[1,m-1]$.  A child of type $i$ has
width $t_i f$, where $f=\lfloor e/m\rfloor$, and operates on exactly
$1/W$ of the parent's logical volume.  Repetitions in the list represent
separate calls.  An ordinary rank-one call has $t_i=1$; a batched block
has the corresponding larger value of $t_i$.  Full-rank edges may retain
their original $m$ rank-one calls, so no same-width recursive call is
needed.  Put
\[
 \Psi(x)=\frac1W\sum_{i=1}^J(t_i/m)^x,
 \qquad t_* = \max_i t_i <m.
\]
If $\Psi(\tau)<1$ for some $0<\tau<1$, then the normalized recurrence
\[
 F(e)\le \frac1W\sum_{i=1}^J F(t_i\lfloor e/m\rfloor)+C,
 \qquad F(e)=O(1)\quad(1\le e<m),
\]
has $F(e)=O(e^\tau)$ for every positive integer $e$.  Indeed, substituting
$F(u)\le A u^\tau$ into an internal call gives at most
$A\Psi(\tau)e^\tau+C$.  Choose $A$ to cover the finite base range and
$C/(1-\Psi(\tau))$.  Floors can only decrease the right-hand side.

\paragraph{A finite integer depth bound.}
Define
\[
 D(e)=0\quad(e<m),\qquad
 D(e)=1+D(t_*\lfloor e/m\rfloor)\quad(e\ge m).
\]
This is an integer algorithm: repeatedly replace $e$ by
$t_*\lfloor e/m\rfloor$ until it is below $m$.  The function $D$ is
nondecreasing, every child has depth bound at most $D(e)-1$, and
$D(e)=O(\log(2e))$, because each internal replacement is at most
$(t_*/m)e$.  A row count divisible by $W^{D(e)}$ is therefore sufficient
for every root-to-leaf path.  At each invocation write the row index as
$Wg+w$ and split by $w$.  Every child has complete within-row ranges,
logical volume exactly $V/W$, and a remaining row count divisible by
$W^{D(e)-1}$.  No padding occurs inside this recursion.

\paragraph{Bit interchange with a remainder.}
At an integer-width call write $e=mf+r$, $0\le r<m$.  If $f=0$, interchange
the bounded number of digits individually.  Otherwise separate the high
$r$ digits of each of the two target chunks.  With $G$ denoting all
intervening spectators, the layout is
\[
 (H^+,H^-,G,D^+,D^-).
\]
First exchange corresponding digits of $H^+,D^+$ and move $H^+$ just
before $H^-$.  The layout becomes
$(D^+,H^+,H^-,G,D^-)$.  Apply the finite circuit to the two length-$mf$
chunks $H^-,D^-$.  Finally move $H^+$ back immediately before the final
$H^-$, obtaining
\[
 (D^+,D^-,G,H^+,H^-).
\]
All other fields are preserved.  These moves use fewer than a fixed
multiple of $m$ single-digit operations and cost $O(V)$.  The main chunks
split into $m$ equal width-$f$ fields, so every rational shear and every
batched contiguous-block shear has exactly the required field shape.
A child retains the unused main digits and all remainders as spectators.
Thus the preceding moment recurrence applies literally at every integer
width; there is no per-level numerical padding of either chunk range.

\paragraph{Supplying the bit rows from existing digits.}
For an arbitrary width $e$, compute $D(e)$ and choose the least nonnegative
integer $\rho$ such that
\[
 q^{2\rho}\ge W^{D(e)}.
\]
Both numbers can be constructed by fixed-radix integer arithmetic.  Since
$D(e)=O(\log e)$, we have $\rho=O(\log e)$.  For sufficiently large $e$,
$2\rho<e$.  The exceptional set of smaller widths is fixed and finite;
individual-digit interchange there is absorbed by the asymptotic constant.

Exchange the high $\rho$ digits of the two chunks, then gather their two
high parts into one row field before both remaining chunks, exactly as in
the preceding layout operation.  This row field has range $q^{2\rho}$.
Within each original prefix, pad its range to
\[
 R'=W^{D(e)}\left\lceil q^{2\rho}/W^{D(e)}\right\rceil
       <2q^{2\rho}.
\]
Every added row includes the entire remaining address set.  Recurse once
at the remaining width $e-\rho$, which has depth bound at most $D(e)$;
then remove the added rows and undo the gathering move.  Each completed
invocation preserves its row index, so padded zero rows are zero again
when they are removed.

The possibly large exponent in $W^{D(e)}=e^{O(1)}$ does not multiply the
logical volume by that polynomial.  The row digits were already part of
the two input chunks.  Only rounding the row count adds data, by a factor
strictly below two, once at the outer boundary.  Exchanging and gathering
$O(\log e)$ high digits costs $O(V\log(2e))$, which is
$O(Ve^\tau)$.  The radix conversion from binary chunk ranges to the fixed
radix $q$ retains the original constant factor, less than $q^2$.
Consequently the batched construction is an arbitrary-width interchange
procedure with the asserted exponent $\tau$ and no additional power of
its volume or width.

\paragraph{Fixed tapes and descriptors.}
Use the original depth-first scheduler: park only the other role streams
at the active stack top, run the child on one complete role stream, and
restore the parked streams on return.  The number of child types is
fixed, so parking, copying, cleanup and the finite arithmetic transforms
contribute $O(V)$ per node.  Before descending, retain the two current
chunks and merge consecutive spectator fields; the descriptor has a fixed
number of field intervals.  Its integer lengths have $O(\log V_0)$ bits.
All descendants retain the complete ranges of both root main chunks,
even when most digits have become spectators.  Their volume is therefore
at least $q^{2(e-\rho)}$ in the outer construction, while
$j\log W=O(\log e)$ at depth $j$.  Hence
$\log V_0=\log V_j+j\log W=O(\log(2V_j))$.
Polynomial descriptor work is absorbed by $O(V_j)$, as in the original
fixed-tape proof.  The larger fixed constant in the depth bound creates
no new tape and no growing alphabet.
